Recall from Complements~\ref{sga3:XXI:2.3.7} that the bilinear form on
\(\VR(R)\) defined by
\[
 \langle x,y\rangle=(p(x),y)
\]
is a Euclidean inner product. Moreover, for \(\alpha\in R\),
\[
 \langle\alpha,y\rangle
 =\frac{\ell(\alpha)}{2}(\alpha^\ast,y),
\]
so \(\langle\alpha,y\rangle\) and \((\alpha^\ast,y)\) have the same
sign.

\begin{lemma}\label{sga3:XXI:3.2.11}
If \(\alpha,\beta\in\Delta\), then
\[
 (\beta^\ast,\alpha)\leq0
 \qquad\text{and hence}\qquad
 \langle\beta,\alpha\rangle\leq0.
\]
\end{lemma}

\begin{proof}
Indeed,
\[
 s_\beta(\alpha)
 =\alpha-(\beta^\ast,\alpha)\beta\in R,
\]
and therefore \((\beta^\ast,\alpha)\leq0\) by
Lemma~\ref{sga3:XXI:3.2.10}.
\end{proof}

We now show that \(\Delta\) is linearly independent. Otherwise,
Lemma~\ref{sga3:XXI:3.2.11} would imply, exactly as in the proof of
Proposition~\ref{sga3:XXI:3.1.5}, the existence of a nontrivial
relation
\[
 \sum_{\alpha\in\Delta}m(\alpha)\alpha=0,
 \qquad m(\alpha)\in\mathbb N.
\]
If \(m(\alpha_0)\geq1\), then
\[
 -\alpha_0
 =\bigl(m(\alpha_0)-1\bigr)\alpha_0
  +\sum_{\alpha\neq\alpha_0}m(\alpha)\alpha.
\]
By (P 3), this gives
\(\alpha_0\in R_+\cap(-R_+)\), contradicting (P 1).
Thus \(\Delta\) is a base of \(\cR\), completing the proof of
Theorem~\ref{sga3:XXI:3.2.8}.

\begin{corollary}\label{sga3:XXI:3.2.12}
Let \(R_+\) be a system of positive roots, let \(\Delta\) be a base of
\(\cR\), and let \(w\in W\). Then
\[
 \Delta\subset R_+
 \quad\Longleftrightarrow\quad
 R_+=\mathscr{P}(\Delta)
 \quad\Longleftrightarrow\quad
 \Delta=\mathscr{S}(R_+),
\]
\[
 \mathscr{P}\bigl(\operatorname{ind}(\Delta^\ast)\bigr)
   =\mathscr{P}(\Delta)^\ast,
 \qquad
 \mathscr{S}(R^\ast_+)
   =\operatorname{ind}\bigl(\mathscr{S}(R_+)^\ast\bigr),
\]
and
\[
 \mathscr{S}\bigl(w(R_+)\bigr)=w\bigl(\mathscr{S}(R_+)\bigr),
 \qquad
 \mathscr{P}\bigl(w(\Delta)\bigr)=w\bigl(\mathscr{P}(\Delta)\bigr).
\]
\end{corollary}

\begin{definition}\label{sga3:XXI:3.2.13}
Once a system of simple roots \(\Delta\) has been chosen, the elements
of \(\mathscr{P}(\Delta)\) are called \emph{positive}. Once a system
of positive roots \(R_+\) has been chosen, the elements of
\(\mathscr{S}(R_+)\) are called \emph{simple}.
\end{definition}

\begin{corollary}\label{sga3:XXI:3.2.14}
Let \(R_+\) be a system of positive roots, and let \(\alpha\in R_+\).
The following conditions are equivalent:
\begin{enumerate}
\item[\SGAInlineRefTarget{sga3:XXI:3.2.14:item:i}(i)] \(\alpha\) is simple, that is,
  \(\alpha\in\mathscr{S}(R_+)\);
\item[\SGAInlineRefTarget{sga3:XXI:3.2.14:item:ii}(ii)] \(\alpha\) is not the sum of two elements of \(R_+\);
\item[\SGAInlineRefTarget{sga3:XXI:3.2.14:item:iii}(iii)] \(R_+\setminus\{\alpha\}\) is closed.
\end{enumerate}
\end{corollary}

\begin{proof}
The equivalence of \SGARefLink{sga3:XXI:3.2.14:item:i}{(i)} and \SGARefLink{sga3:XXI:3.2.14:item:ii}{(ii)} follows immediately from
Lemma~\ref{sga3:XXI:3.2.9}; the equivalence of \SGARefLink{sga3:XXI:3.2.14:item:ii}{(ii)} and \SGARefLink{sga3:XXI:3.2.14:item:iii}{(iii)} is
immediate from Definition~\ref{sga3:XXI:3.1.4}.
\end{proof}

\begin{definition}\label{sga3:XXI:3.2.15}
Let \(\Delta\) be a system of simple roots. The sum of the
coefficients in the expression of a root \(\alpha\) relative to
\(\Delta\) is called the \emph{order of \(\alpha\) relative to
\(\Delta\)} and is denoted \(\operatorname{ord}_\Delta(\alpha)\).
\end{definition}

One has
\[
 \alpha\in\mathscr{P}(\Delta)
 \quad\Longleftrightarrow\quad
 \operatorname{ord}_\Delta(\alpha)>0,
 \qquad
 \alpha\in\Delta
 \quad\Longleftrightarrow\quad
 \operatorname{ord}_\Delta(\alpha)=1.
\]

\begin{lemma}\label{sga3:XXI:3.2.16}
Let \(\Delta\) be a system of simple roots and
\(\alpha\in\mathscr{P}(\Delta)\). There is a sequence
\(\alpha_i\in\Delta\), \(i=1,\ldots,m\), such that
\[
 \alpha_1+\alpha_2+\cdots+\alpha_p\in\mathscr{P}(\Delta)
 \quad\text{for }p=1,\ldots,m,
\]
and
\[
 \alpha_1+\alpha_2+\cdots+\alpha_m=\alpha.
\]
Moreover, for every sequence satisfying these conditions,
\[
 m=\operatorname{ord}_\Delta(\alpha).
\]
\end{lemma}

\begin{proof}
This is immediate from Proposition~\ref{sga3:XXI:3.1.2}.
\end{proof}
